\documentclass[10.5pt,a4paper]{article}
\usepackage[T1]{fontenc}
\usepackage[utf8]{inputenc}
\usepackage{mathptmx} % Times-like serif (Elsevier body font)
\usepackage[margin=1.55cm]{geometry}
\usepackage{xcolor}
\usepackage{enumitem}
\usepackage{titlesec}
\usepackage[protrusion=true,expansion=false]{microtype}
\usepackage[official]{eurosym}
\usepackage[hidelinks]{hyperref}

\definecolor{accent}{HTML}{000000}
\definecolor{accent2}{HTML}{000000}
\definecolor{grey}{HTML}{000000}
\hypersetup{colorlinks=true, urlcolor=accent2, linkcolor=accent2}
\hyphenpenalty=10000 \exhyphenpenalty=10000 \tolerance=9999 \emergencystretch=3em
\setlength{\parindent}{0pt}

\titleformat{\section}{\normalsize\bfseries\color{accent}}{}{0em}{}[{\vspace{-0.3ex}\color{accent}\titlerule[0.8pt]}]
\titlespacing*{\section}{0pt}{1.1ex}{0.5ex}
\setlist[itemize]{leftmargin=1.25em, itemsep=0.5pt, topsep=1pt, parsep=0pt, label=\textbullet}

\newcommand{\dateb}[1]{{\small\itshape\color{grey}#1}}
\newcommand{\cvrow}[2]{%
  \noindent\makebox[2.6cm][l]{\dateb{#1}}%
  \parbox[t]{\dimexpr\linewidth-2.6cm\relax}{#2}\par\smallskip}
\pagestyle{empty}

\begin{document}

{\LARGE\bfseries\color{accent} Dr. NKODIA Hardy Medry Dieu-Veil}\\[2pt]
{\itshape\color{accent2} Géologue — Tectonique, géodynamique et géologie structurale}\\[3pt]
{\small\color{grey} Brazzaville, République du Congo \ \textbullet\ \href{mailto:nkodiahardy@gmail.com}{nkodiahardy@gmail.com}}\\[1pt]
{\small \href{https://orcid.org/0000-0001-6919-6863}{ORCID 0000-0001-6919-6863} \ \textbullet\ \href{https://scholar.google.com/citations?user=gYOI-FkAAAAJ&hl=fr}{Google Scholar} \ \textbullet\ \href{https://www.researchgate.net/profile/Hardy-Medry-Dieu-Veil-Nkodia}{ResearchGate} \ \textbullet\ \href{https://nkodiahardy.com}{nkodiahardy.com}}

\section{Profil}
Enseignant-chercheur en géologie structurale et tectonique : analyse des paléocontraintes, réactivation des failles et aléa sismique intraplaque en Afrique centrale. \textit{Axes} : réservoirs failles/fractures (Inkisi, Pool)~; géomorphologie du bassin du Congo~; minéralisations et aquifères~; aléa sismique de la marge riftée d'Afrique centrale.

\section{Formation}
\cvrow{2019 – 2024}{\textbf{Doctorat en Géosciences} — Université Marien Ngouabi (Congo) \& Musée royal de l'Afrique centrale (Belgique). Mention très honorable avec félicitations du jury.}
\cvrow{2015 – 2017}{\textbf{Master en Géosciences} — Université Marien Ngouabi. Mention Bien~; major de promotion.}
\cvrow{2012 – 2015}{\textbf{Licence en Géosciences} — Université Marien Ngouabi. Mention Assez bien~; 2\textsuperscript{e} de promotion.}
\cvrow{2011}{\textbf{Baccalauréat série D} — Lycée Chaminade, Brazzaville. Mention Bien.}

\section{Expérience académique}
\cvrow{2025 – présent}{\textbf{Assistant} (tenure-track, en voie d'Assistant Professor) — Université Marien Ngouabi, FST, Dép. de Géologie (Congo). Enseignement et recherche en géologie structurale~; encadrement de Master.}
\cvrow{2025 – 2026}{\textbf{Chercheur post-doctoral} — Pukyong National University, Busan (Corée du Sud), en congé de recherche de l'UMNG. Champ de contraintes et structures sismogènes de la péninsule coréenne.}
\cvrow{Août – sept. 2026}{\textbf{Séjour de recherche} — MRAC, Tervuren (Belgique). Chaîne Ouest-congolienne et karst.}
\cvrow{Févr.–mars 2025}{\textbf{Chercheur post-doctoral} (séjour) — Musée royal de l'Afrique centrale (MRAC), Tervuren (Belgique).}
\cvrow{2018 – 2024}{\textbf{Enseignant-chercheur} — Université Marien Ngouabi (Congo). Cours de Géologie structurale (L2)~; co-organisation d'une conférence internationale.}
\cvrow{2019 – 2021}{\textbf{Assistant de recherche} — MRAC, Tervuren (Belgique). Gestion d'un budget de \euro{}15\,000 sur 3 ans~; collaborations internationales.}
\cvrow{2019 – 2021}{\textbf{Enseignant} — École Africaine de Développement (EAD), Brazzaville. Cours de géologie structurale.}

\section{Financements \& distinctions}
\begin{itemize}
\item \textbf{AMORCE-GEO (2025–2027)} — projet financé par l'ARES (Belgique)~; rôle : co-organisateur et formateur.
\item Doctorat avec félicitations du jury (2024)~; major de promotion (Master, 2017).
\end{itemize}

\section{Publications — sélection (articles à comité de lecture)}
{\small\color{grey}Liste complète : 12 articles à comité de lecture (voir ORCID / Google Scholar).}
\begin{enumerate}[leftmargin=1.5em, itemsep=2pt, topsep=2pt, label=\arabic*.]
\item \textbf{Nkodia, H.\,M.\,D.-V.}, Hategekimana, F., Naik, S.\,P., Cheon, Y., Peace, A.\,L., Bae, S.-Y., Kandregula, R.\,S., \& Kim, Y.-S. (2026). Vertical kinematic partitioning in intraplate fault systems: evidence from the SE Korean Peninsula. \emph{Journal of Structural Geology}, \textbf{211}, 105778. \href{https://doi.org/10.1016/j.jsg.2026.105778}{doi:10.1016/j.jsg.2026.105778}
\item \textbf{Nkodia, H.\,M.\,D.-V.}, Park, K., Naik, S.\,P., Peace, A.\,L., \& Kim, Y.-S. (2026). Assessing the contemporary stress field and seismogenic structures of the Korean peninsula using focal-mechanism inversion and slip-tendency analysis. \emph{Tectonophysics}, \textbf{934}, 231261. \href{https://doi.org/10.1016/j.tecto.2026.231261}{doi:10.1016/j.tecto.2026.231261}
\item Colet, M., Kolawole, F., Ajala, R., Delvaux, D., \& \textbf{Nkodia, H.\,M.\,D.-V.} (2025). Active crustal deformation across a nucleating extensional microplate, D.\,R. Congo, East Africa. \emph{Tectonics}, 44, e2025TC008815. \href{https://doi.org/10.1029/2025TC008815}{doi:10.1029/2025TC008815}
\item \textbf{Nkodia, H.\,M.\,D.-V.}, Boudzoumou, F., Miyouna, T., \emph{et al.} (2024). Brittle faulting and tectonic stress history on the western margin of the Congo Basin between Kinshasa and Brazzaville. \emph{Tectonophysics}, 877, 230282. \href{https://doi.org/10.1016/j.tecto.2024.230282}{doi:10.1016/j.tecto.2024.230282}
\item Mavoungou, Y.\,B., \textbf{Nkodia, H.\,M.\,D.-V.}, Watha-Ndoudy, N., \& Bolarinwa, A.\,T. (2024). Lineament extraction and paleostress analysis in the Bikélélé iron deposit (Chaillu Massif, Congo). \emph{Modeling Earth Systems and Environment}, 10(2), 1993–2009. \href{https://doi.org/10.1007/s40808-023-01883-3}{doi:10.1007/s40808-023-01883-3}
\item \textbf{Nkodia, H.\,M.\,D.-V.}, Miyouna, T., Kolawole, F., \emph{et al.} (2022). Seismogenic fault reactivation in western Central Africa: insights from regional stress analysis. \emph{Geochemistry, Geophysics, Geosystems}, 23(11), e2022GC010377. \href{https://doi.org/10.1029/2022GC010377}{doi:10.1029/2022GC010377}
\item \textbf{Nkodia, H.\,M.\,D.-V.}, Miyouna, T., Delvaux, D., \& Boudzoumou, F. (2020). Flower structures in sandstones of the Palaeozoic Inkisi Group (Brazzaville, Congo). \emph{South African Journal of Geology}, 123(4), 531–550. \href{https://doi.org/10.25131/sajg.123.0038}{doi:10.25131/sajg.123.0038}
\item Miyouna, T., \textbf{Nkodia, H.\,M.\,D.-V.}, Essouli, O.\,F., \emph{et al.} (2018). Strike-slip deformation in the Inkisi Formation, Brazzaville, Republic of Congo. \emph{Cogent Geoscience}, 4(1), 1542762. \href{https://doi.org/10.1080/23312041.2018.1542762}{doi:10.1080/23312041.2018.1542762}
\end{enumerate}

\section{Conférence invitée}
\begin{itemize}
\item \textbf{6 juillet 2026} — « Ocean–continent tectonic stresses and earthquake hazards of the Central African rifted margin ». Séminaire en ligne \emph{Earthquake Hazards of Rifted Margins}, Columbia University / Lamont-Doherty Earth Observatory. \href{https://www.youtube.com/watch?v=lnRf0Jz0QbM}{Enregistrement}.
\end{itemize}

\section{Compétences, langues \& responsabilités}
\begin{itemize}
\item Analyse tectonique/structurale et paléocontraintes (slip tendency)~; télédétection et SIG (DEM ALOS-PALSAR, imagerie satellitaire/drone)~; pétrologie et géologie de terrain~; Python (figures publiables), Adobe Illustrator.
\item \textit{Certifications} : Gestion de projet (Google, 2024)~; rédaction/publication scientifique (École Polytechnique Paris, 2023)~; télédétection radar (IGNFI/AFD, 2019).
\item \textit{Langues} : français (langue maternelle), anglais (professionnel). \quad \textit{Responsabilités} : fondateur \& président de Kongo Science~; fondateur du LGRN (UMNG).
\end{itemize}

\section{Références}
\begin{itemize}
\item \textbf{Pr. F. Boudzoumou} (UMNG).
\item \textbf{Dr. D. Delvaux} — Editor-in-Chief, \emph{J. African Earth Sciences}~; MRAC, Tervuren.
\item \emph{Coordonnées des référents disponibles sur demande.}
\end{itemize}

\end{document}
